\documentclass[11pt,a4paper]{article}
\usepackage[top=1.6cm,bottom=1.4cm,left=1.3cm,right=1.3cm,headsep=8pt,footskip=16pt]{geometry}
\usepackage{enumitem}
\usepackage{amssymb}
\usepackage{array}
\usepackage[table]{xcolor}
\usepackage{fancyhdr}
\usepackage{multicol}
\usepackage{polyglossia}
\setdefaultlanguage[calendar=gregorian, numerals=maghrib]{arabic}
\setotherlanguage{english}
\IfFontExistsTF{Amiri}{%
  \newfontfamily\arabicfont[Script=Arabic, Scale=1.08]{Amiri}
}{%
  \newfontfamily\arabicfont[Script=Arabic, Scale=1.08, Path=../fonts/, Extension=.ttf,
    UprightFont=*-Regular, BoldFont=*-Bold, ItalicFont=*-Italic,
    BoldItalicFont=*-BoldItalic]{Amiri}
}
\IfFontExistsTF{AlBattar}{%
  \newfontfamily\khatt[Script=Arabic, Scale=1.15]{AlBattar}
}{%
  \newfontfamily\khatt[Script=Arabic, Scale=1.15, Path=../fonts/,
    BoldFont=ae_AlBattar.ttf]{ae_AlBattar.ttf}
}
\renewcommand{\headrulewidth}{0.8pt}
\renewcommand{\footrulewidth}{0.8pt}
\fancyhf{}
\rhead{\khatt المدرسة العليا للأساتذة ببوزريعة}
\lhead{\khatt ملحقة بومرداس}
\lfoot{\khatt الأستاذ: د. بن عبد الرزاق يوسف}
\cfoot{\thepage/2}
\pagestyle{fancy}
\parindent 0pt
\setlength{\columnsep}{0.6cm}
\setlength{\columnseprule}{0.4pt}
\setlist[enumerate,1]{label=\textbf{\arabic*.}, leftmargin=*, itemsep=5pt, topsep=2pt}
\begin{document}

\begin{center}
{\Large\khatt\textbf{الفرض الخامس في الذكاء الاصطناعي}}\\[2pt]
{\normalsize أسئلة متعددة الاختيارات ــ الفصول من 1 إلى 6 ــ 15 سؤالًا}
\end{center}
\vspace{-4pt}\rule{\textwidth}{0.6pt}\\
{\small\textbf{التعليمة:} لكل سؤال إجابة صحيحة واحدة فقط من بين خمسة اختيارات.
اكتب \textbf{حرف} الإجابة الصحيحة في جدول الإجابات أسفل الصفحة الثانية.
لكل سؤال نقطة واحدة (المجموع 15 نقطة).}
\vspace{-2pt}
\begin{multicols}{2}

\noindent\begin{minipage}{\linewidth}
\hangindent=0.7cm\hangafter=1\makebox[0.7cm][r]{\textbf{1.}}\textbf{أي مما يلي مصدر موثوق للتحقق من إعراب كلمة؟}
\begin{itemize}[leftmargin=0.95cm, labelsep=4pt, itemsep=0pt, topsep=2pt, parsep=0pt]
  \item[\textbf{أ)}] كتاب نحو معتمد ومعروف
  \item[\textbf{ب)}] إجابة أداة ذكاء اصطناعي ثانية
  \item[\textbf{ج)}] منشور على شبكة اجتماعية
  \item[\textbf{د)}] مدونة مجهولة الكاتب
  \item[\textbf{هـ)}] تعليق قارئ في منتدى
\end{itemize}
\end{minipage}\par\vspace{7pt plus 4pt}
\noindent\begin{minipage}{\linewidth}
\hangindent=0.7cm\hangafter=1\makebox[0.7cm][r]{\textbf{2.}}\textbf{عبارة «في 150 كلمة، بلغة بسيطة» تمثّل عنصر:}
\begin{itemize}[leftmargin=0.95cm, labelsep=4pt, itemsep=0pt, topsep=2pt, parsep=0pt]
  \item[\textbf{أ)}] السياق
  \item[\textbf{ب)}] التنسيق
  \item[\textbf{ج)}] الدور
  \item[\textbf{د)}] المهمة
  \item[\textbf{هـ)}] القيود
\end{itemize}
\end{minipage}\par\vspace{7pt plus 4pt}
\noindent\begin{minipage}{\linewidth}
\hangindent=0.7cm\hangafter=1\makebox[0.7cm][r]{\textbf{3.}}\textbf{بماذا يبدأ البحث الجيد عن المعلومة؟}
\begin{itemize}[leftmargin=0.95cm, labelsep=4pt, itemsep=0pt, topsep=2pt, parsep=0pt]
  \item[\textbf{أ)}] بسؤال واضح
  \item[\textbf{ب)}] بالتلخيص
  \item[\textbf{ج)}] بمحرك البحث
  \item[\textbf{د)}] بنسخ المقالات
  \item[\textbf{هـ)}] بالنتيجة الأولى
\end{itemize}
\end{minipage}\par\vspace{7pt plus 4pt}
\noindent\begin{minipage}{\linewidth}
\hangindent=0.7cm\hangafter=1\makebox[0.7cm][r]{\textbf{4.}}\textbf{أي استعمال من الاستعمالات الآتية ذكره الدرس لأستاذ اللغة العربية؟}
\begin{itemize}[leftmargin=0.95cm, labelsep=4pt, itemsep=0pt, topsep=2pt, parsep=0pt]
  \item[\textbf{أ)}] وضع العلامات النهائية للتلاميذ آليًا
  \item[\textbf{ب)}] تعويض الأستاذ في القسم تعويضًا كاملًا
  \item[\textbf{ج)}] اقتراح أمثلة جديدة لقاعدة نحوية
  \item[\textbf{د)}] اختيار المنهاج الدراسي بدل الوزارة
  \item[\textbf{هـ)}] تصحيح أوراق الامتحان دون قراءتها
\end{itemize}
\end{minipage}\par\vspace{7pt plus 4pt}
\noindent\begin{minipage}{\linewidth}
\hangindent=0.7cm\hangafter=1\makebox[0.7cm][r]{\textbf{5.}}\textbf{نموذج تصحيح تدرّب على نصوص صحفية حديثة فقط، فعدّ كلمات من الشعر الجاهلي أخطاءً. ما السبب؟}
\begin{itemize}[leftmargin=0.95cm, labelsep=4pt, itemsep=0pt, topsep=2pt, parsep=0pt]
  \item[\textbf{أ)}] النموذج لا يستطيع قراءة النصوص العربية
  \item[\textbf{ب)}] الحاسوب كان بطيئًا أثناء مرحلة الاختبار
  \item[\textbf{ج)}] لم يرَ هذا النوع من النصوص أثناء التدريب
  \item[\textbf{د)}] بيانات الاختبار كانت كبيرة جدًا ومتنوعة
  \item[\textbf{هـ)}] الشعر الجاهلي مليء بالأخطاء اللغوية
\end{itemize}
\end{minipage}\par\vspace{7pt plus 4pt}
\noindent\begin{minipage}{\linewidth}
\hangindent=0.7cm\hangafter=1\makebox[0.7cm][r]{\textbf{6.}}\textbf{أي أمر من الأوامر الآتية يحتوي على عنصر التنسيق؟}
\begin{itemize}[leftmargin=0.95cm, labelsep=4pt, itemsep=0pt, topsep=2pt, parsep=0pt]
  \item[\textbf{أ)}] «في 100 كلمة على الأكثر»
  \item[\textbf{ب)}] «نظّم الإجابة في جدول»
  \item[\textbf{ج)}] «تلاميذي في السنة الثالثة»
  \item[\textbf{د)}] «تصرّف كمفتش تربوي خبير»
  \item[\textbf{هـ)}] «اشرح درس الحال للتلاميذ»
\end{itemize}
\end{minipage}\par\vspace{7pt plus 4pt}
\noindent\begin{minipage}{\linewidth}
\hangindent=0.7cm\hangafter=1\makebox[0.7cm][r]{\textbf{7.}}\textbf{ما القاعدة البسيطة التي يقترح الدرس أن نعلّمها للتلاميذ؟}
\begin{itemize}[leftmargin=0.95cm, labelsep=4pt, itemsep=0pt, topsep=2pt, parsep=0pt]
  \item[\textbf{أ)}] «ما تكتبه الأداة ملكي أقدّمه باسمي»
  \item[\textbf{ب)}] «النتيجة النهائية أهم من طريقة الوصول إليها»
  \item[\textbf{ج)}] «الأداة تساعدني على التفكير ولا تفكر بدلًا مني»
  \item[\textbf{د)}] «الأداة تفكر بدلًا مني فأرتاح من التفكير»
  \item[\textbf{هـ)}] «لا حاجة إلى التصريح باستعمال الأداة»
\end{itemize}
\end{minipage}\par\vspace{7pt plus 4pt}
\noindent\begin{minipage}{\linewidth}
\hangindent=0.7cm\hangafter=1\makebox[0.7cm][r]{\textbf{8.}}\textbf{ما الكلمات المفتاحية المناسبة للسؤال: «كيف يساعد الذكاء الاصطناعي على تحسين التعبير الكتابي؟»}
\begin{itemize}[leftmargin=0.95cm, labelsep=4pt, itemsep=0pt, topsep=2pt, parsep=0pt]
  \item[\textbf{أ)}] تحسين، كيف، هل، يساعد، على، في، من، إلى
  \item[\textbf{ب)}] ما، هل، كيف، لماذا، متى، أين، من، على
  \item[\textbf{ج)}] على، الكتابي، كيف، يساعد، ما، هل، لماذا
  \item[\textbf{د)}] كيف، يساعد، على، تحسين، هل، ما، لماذا، متى
  \item[\textbf{هـ)}] الذكاء الاصطناعي، التعبير الكتابي، \textenglish{Writing skills}
\end{itemize}
\end{minipage}\par\vspace{7pt plus 4pt}
\noindent\begin{minipage}{\linewidth}
\hangindent=0.7cm\hangafter=1\makebox[0.7cm][r]{\textbf{9.}}\textbf{السؤال «هل هدفه الإعلام، أم الإقناع، أم البيع؟» يخص معيار:}
\begin{itemize}[leftmargin=0.95cm, labelsep=4pt, itemsep=0pt, topsep=2pt, parsep=0pt]
  \item[\textbf{أ)}] التاريخ
  \item[\textbf{ب)}] الكاتب
  \item[\textbf{ج)}] الجهة
  \item[\textbf{د)}] الهدف
  \item[\textbf{هـ)}] الدقة
\end{itemize}
\end{minipage}\par\vspace{7pt plus 4pt}
\noindent\begin{minipage}{\linewidth}
\hangindent=0.7cm\hangafter=1\makebox[0.7cm][r]{\textbf{10.}}\textbf{ما الخطأ المحتمل في «الأرقام» عند قراءة الأداة لصورة؟}
\begin{itemize}[leftmargin=0.95cm, labelsep=4pt, itemsep=0pt, topsep=2pt, parsep=0pt]
  \item[\textbf{أ)}] تحوّل الأرقام كلها إلى تسجيل صوتي
  \item[\textbf{ب)}] تحذف الصورة الأصلية من جهاز المستخدم
  \item[\textbf{ج)}] ترفض قراءة أي رقم موجود في الصورة
  \item[\textbf{د)}] تقرأ رقمًا خطأ، أو تخلط بين عمودين
  \item[\textbf{هـ)}] تغيّر لغة الجهاز إلى اللغة الإنجليزية
\end{itemize}
\end{minipage}\par\vspace{7pt plus 4pt}
\noindent\begin{minipage}{\linewidth}
\hangindent=0.7cm\hangafter=1\makebox[0.7cm][r]{\textbf{11.}}\textbf{ما الترتيب الصحيح لخطوات التحقق من إجابة الأداة؟}
\begin{itemize}[leftmargin=0.95cm, labelsep=4pt, itemsep=0pt, topsep=2pt, parsep=0pt]
  \item[\textbf{أ)}] اقرأ $\leftarrow$ حدّد المعلومة $\leftarrow$ ابحث في مصدر موثوق $\leftarrow$ قارن $\leftarrow$ قرّر
  \item[\textbf{ب)}] حدّد المعلومة $\leftarrow$ اقرأ $\leftarrow$ قارن $\leftarrow$ قرّر $\leftarrow$ ابحث في مصدر موثوق
  \item[\textbf{ج)}] اقرأ $\leftarrow$ ابحث في مصدر موثوق $\leftarrow$ قرّر $\leftarrow$ حدّد المعلومة $\leftarrow$ قارن
  \item[\textbf{د)}] ابحث في مصدر موثوق $\leftarrow$ قارن $\leftarrow$ اقرأ $\leftarrow$ حدّد المعلومة $\leftarrow$ قرّر
  \item[\textbf{هـ)}] اقرأ $\leftarrow$ قرّر $\leftarrow$ حدّد المعلومة $\leftarrow$ قارن $\leftarrow$ ابحث في مصدر موثوق
\end{itemize}
\end{minipage}\par\vspace{7pt plus 4pt}
\noindent\begin{minipage}{\linewidth}
\hangindent=0.7cm\hangafter=1\makebox[0.7cm][r]{\textbf{12.}}\textbf{نموذج حقق نتائج ممتازة على بيانات التدريب، لكنه فشل مع الأمثلة الجديدة. ما التفسير الأقرب؟}
\begin{itemize}[leftmargin=0.95cm, labelsep=4pt, itemsep=0pt, topsep=2pt, parsep=0pt]
  \item[\textbf{أ)}] بيانات الاختبار كانت متنوعة، وهذا عيب
  \item[\textbf{ب)}] النتيجة جيدة، لأن نتائج التدريب هي الأهم
  \item[\textbf{ج)}] كان يجب اختباره ببيانات التدريب نفسها
  \item[\textbf{د)}] حفظ الأمثلة دون أن يكتشف أنماطًا عامة
  \item[\textbf{هـ)}] الفشل سببه أن مراحل التعلم ست لا خمس
\end{itemize}
\end{minipage}\par\vspace{7pt plus 4pt}
\noindent\begin{minipage}{\linewidth}
\hangindent=0.7cm\hangafter=1\makebox[0.7cm][r]{\textbf{13.}}\textbf{أي أمر يحتوي على العناصر الخمسة كلها؟}
\begin{itemize}[leftmargin=0.95cm, labelsep=4pt, itemsep=0pt, topsep=2pt, parsep=0pt]
  \item[\textbf{أ)}] «تصرّف كأستاذ نحو، ونظّم شرحك في جدول من عمودين، مع أمثلة من الحياة»
  \item[\textbf{ب)}] «تصرّف كأستاذ نحو. اشرح المفعول المطلق لتلاميذ الرابعة، في 120 كلمة، في جدول»
  \item[\textbf{ج)}] «تصرّف كأستاذ نحو. تلاميذي في السنة الرابعة. في 120 كلمة، وفي جدول»
  \item[\textbf{د)}] «اشرح المفعول المطلق لتلاميذ السنة الرابعة متوسط، في 120 كلمة، بلغة بسيطة»
  \item[\textbf{هـ)}] «تصرّف كأستاذ نحو. اشرح المفعول المطلق بلغة بسيطة، مع جدول للأمثلة»
\end{itemize}
\end{minipage}\par\vspace{7pt plus 4pt}
\noindent\begin{minipage}{\linewidth}
\hangindent=0.7cm\hangafter=1\makebox[0.7cm][r]{\textbf{14.}}\textbf{أستاذ يريد تعليمة لواجب منزلي تعلّم التلاميذ الاستعمال الأخلاقي. أي تعليمة أنسب؟}
\begin{itemize}[leftmargin=0.95cm, labelsep=4pt, itemsep=0pt, topsep=2pt, parsep=0pt]
  \item[\textbf{أ)}] «استعملوا الأداة لكل شيء، فالمهم هو النتيجة النهائية»
  \item[\textbf{ب)}] «استعينوا بالأداة للفهم، ثم اكتبوا بأسلوبكم وصرّحوا»
  \item[\textbf{ج)}] «لا تصرّحوا بالأداة أبدًا لأن ذلك ينقص من علامتكم»
  \item[\textbf{د)}] «يكفي ذكر اسم الأداة دون مراجعة أي شيء مما كتبته»
  \item[\textbf{هـ)}] «انسخوا إجابة الأداة ثم غيّروا بعض الكلمات فيها»
\end{itemize}
\end{minipage}\par\vspace{7pt plus 4pt}
\noindent\begin{minipage}{\linewidth}
\hangindent=0.7cm\hangafter=1\makebox[0.7cm][r]{\textbf{15.}}\textbf{أي حالة تمثل خطأً في «العلاقات» عند قراءة الأداة لمخطط بياني؟}
\begin{itemize}[leftmargin=0.95cm, labelsep=4pt, itemsep=0pt, topsep=2pt, parsep=0pt]
  \item[\textbf{أ)}] قرأت السنة 1962 على أنها 1926 في المحور
  \item[\textbf{ب)}] قرأت الرقم 35 على أنه 53 في أحد الأعمدة
  \item[\textbf{ج)}] خلطت بين قيمتي عمودين متجاورين في المخطط
  \item[\textbf{د)}] قرأت حرف الفاء قافًا في عنوان المخطط
  \item[\textbf{هـ)}] ادّعت أن ارتفاع العدد سبّب انخفاض النتائج
\end{itemize}
\end{minipage}\par\vspace{7pt plus 4pt}
\end{multicols}
\vfill
\begin{center}
{\khatt\textbf{جدول الإجابات}}\\[3pt]
\renewcommand{\arraystretch}{1.6}
\begin{tabular}{|>{\centering\arraybackslash}m{1.9cm}|*{15}{>{\centering\arraybackslash}m{0.62cm}|}}
\hline
\textbf{رقم السؤال} & \textbf{1} & \textbf{2} & \textbf{3} & \textbf{4} & \textbf{5} & \textbf{6} & \textbf{7} & \textbf{8} & \textbf{9} & \textbf{10} & \textbf{11} & \textbf{12} & \textbf{13} & \textbf{14} & \textbf{15} \\ \hline
\textbf{الإجابة} &  &  &  &  &  &  &  &  &  &  &  &  &  &  &  \\[8pt] \hline
\end{tabular}

\vspace{8pt}
\renewcommand{\arraystretch}{2.2}
\begin{tabular}{|>{\centering\arraybackslash}m{7.5cm}|>{\centering\arraybackslash}m{3.2cm}|>{\centering\arraybackslash}m{3.2cm}|}
\hline
\textbf{الاسم واللقب} & \textbf{الفوج} & \textbf{العلامة} \\ \hline
 & & \hfill /15 \\ \hline
\end{tabular}
\end{center}
\end{document}
